\scanpage{89}
\uline{Cor} Si de plus les $b_*$ existent \inserted{($b\in\mathrm{Ob}\,B$)}
[p.ex. si $p$ est aussi cofibrant, et les ``petites'' sommes
dans les $\unsure{E|\mathcal E}_b$ existent), alors \unsure{le|les} ``foncteur
\unsure{fibre|fibres}'' $\underleftarrow{\mathrm{Lim}}\ E/B\xrightarrow{b_c^*}E_b$ a un adj. à
gauche $b_*^c$.

\uline{Cor} Supposons que $\underline{\Gamma}$ ait une \deleted{petite} famille
de générateurs \inserted{stricts} $(X_\lambda)$, alors les $\unsure{i_*|i}(X_\lambda)$
forment une famille de générateurs \inserted{stricts} de
$\underleftarrow{\mathrm{Lim}}$. En particulier, si $p$ cofibrant, et si
les $E_b$ ont des petites sommes et de
petites familles de générateurs stricts, (de sorte
qu'il en \unsure{soi.|soit|est} de $=$ dans $\underline{\Gamma}$), alors les $=\mathrm{s}$
conditions sont remplies pour $\underleftarrow{\mathrm{Lim}}$, et on trouve
un syst. de générateurs \inserted{stricts} en prenant pour \unsure{tt|tout}
$b\in\mathrm{Ob}\,B$, un ens $\mathrm E_b$ de gén. stricts de $\mathrm E_b$,
et prenant $\displaystyle\bigcup_{b\in\mathrm{Ob}\,B}b_*^c(\mathrm E_b)$.

% archive sheet number: 88
